\section{Related Work}
\label{sec:related_work}

\paragraph{Ambiguity in NLP and requirements engineering.}
Ambiguity has been studied extensively in both computational linguistics and software engineering. Word sense disambiguation \citep{navigli_wsd_2009} addresses lexical ambiguity at the token level; our lexical type extends this to policy-specific polysemy. \citet{li_ambiguity_taxonomy_2024} propose an 11-type NLP taxonomy covering syntactic, semantic, and pragmatic ambiguity, from which we draw our linguistic-layer types. In requirements engineering (RE), \citet{massey_ambiguity_re_2014} classify six ambiguity types in regulatory texts, and \citet{berry_kamsties_2004} introduce the distinction between nocuous ambiguity (where stakeholders diverge on interpretation) and innocuous ambiguity (where they converge). Computational detection of nocuous ambiguity has progressed from heuristic methods \citep{chantree_2006,willis_nocuous_2008} through NLP-based cross-domain detectors \citep{ferrari_ambiguity_2019,ferrari_defects_2018} to transformer-based approaches for anaphoric ambiguity \citep{ezzini_taphsir_2022,ezzini_anaphoric_2022}. \citet{zhao_nlp4re_2021} survey the broader NLP4RE landscape. Our work extends Berry and Kamsties' framework to LLM agents, where the absence of shared organizational context makes previously innocuous ambiguity nocuous.

\paragraph{Psycholinguistics of ambiguity processing.}
\citet{ferreira_good_enough_2002} show that human comprehenders build shallow, ``good-enough'' representations rather than full parses, paralleling our observation that LLM agents adopt surface interpretations. \citet{frazier_rayner_1982} demonstrate early commitment to a single parse; \citet{li_garden_path_llm_2024} provide direct evidence that LLMs exhibit analogous garden-path effects. \citet{liu_ambiguity_2023} demonstrate that language models fail to represent ambiguity, collapsing multiple interpretations into a single output. These findings ground our \emph{Disambiguation Deficit} result: agents process ambiguous input without meta-cognitive awareness but, unlike humans, lack social channels to recover. \citet{suri_structured_2025} show that structured uncertainty estimation improves coverage by 7--39\%, while \citet{sehwag_aspi_2026} demonstrate that clarification-seeking introduces prompt injection vulnerabilities---a tension for disambiguation channel design.

\paragraph{Agent safety evaluation.}
The agent safety evaluation landscape includes benchmarks for risk awareness \citep{rjudge_2024}, tool-use safety \citep{toolemu_2024}, indirect prompt injection \citep{injecagent_2024}, safe task planning \citep{yin_safeagentbench_2025}, and policy-invisible violations \citep{phantompolicy_2026}. \citet{tamkin_evaluating_2023} evaluate how language models navigate policy-governed decisions, \citet{shieldagent_2025} extract formal safety rules from policy documents, and AgentAuditor \citep{agentauditor_2025} achieves human-level safety evaluation. Our work differs by focusing on the \emph{source specification} rather than the agent's execution: we measure how ambiguity in the policy itself causes safety violations, complementing work that assumes clean policies.

\paragraph{Policy specification for agents.}
DRIP-R \citep{drip_r_2026} is the closest concurrent work: it benchmarks 400 policy-ambiguous retail scenarios across five RE ambiguity types, measuring decision quality rather than safety violations. Our work differs in three respects: (1) our matched-pair design establishes causality rather than correlating ambiguity presence with outcomes; (2) we measure safety violations rather than decision quality, enabling failure mechanism analysis; (3) our two-layer taxonomy adds linguistic-layer types, and we quantify the model$\times$type interaction. We map DRIP-R's five RE ambiguity categories to our six-type taxonomy in Appendix~\ref{app:drip_mapping}. PolicyBank \citep{policybank_2026} studies iterative policy refinement, closing 82\% of the gap to a human oracle. AgentPex \citep{agentpex_2026} detects procedural violations in agentic traces, finding that 83\% of perfect-outcome traces contain violations. MANTRA \citep{mantra_2026} synthesizes SMT-validated compliance benchmarks from procedural manuals. $\tau$-bench \citep{tau_bench_2024} and its successor $\tau^2$-bench \citep{tau2_bench_2025} provide the conversational agent evaluation platforms. 